% chapters/08/section-03.tex
\section{Menormalkan Arah dan Membatasi Kecepatan}
\label{sec:menormalkan}

Setelah jarak dan arah diketahui, langkah berikutnya adalah mengubah keduanya
menjadi kecepatan. Di sinilah terletak kesalahan yang paling sering terjadi
pada perilaku mengejar.

\subsection{Cara yang Salah}

Cara yang paling mudah dipikirkan adalah memakai vektor jarak itu langsung
sebagai kecepatan.

\begin{kodekeluaran}
self.velocity = to_player          # cara yang salah
\end{kodekeluaran}

Cara itu tampak masuk akal: vektornya memang menunjuk ke pemain, dan arahnya
benar. Persoalannya terletak pada \textbf{besarnya}. Panjang vektor itu sama
dengan jarak antara NPC dan pemain, sehingga kecepatannya ikut berubah
menurut jarak: NPC bergerak lambat ketika pemain jauh, dan makin cepat ketika
pemain mendekat. Pada jarak 400 piksel, kecepatannya 400 piksel per detik;
pada jarak 75 piksel, kecepatannya 75 piksel per detik.

\subsection{Cara yang Benar}

Cara yang benar adalah memisahkan arah dari besar kecepatan. Arah diambil dari
vektor itu, tetapi besarnya ditentukan oleh nilai kecepatan yang Anda pilih,
bukan oleh jarak.

\begin{kodekeluaran}
self.velocity = to_player.normalize() * self.speed
\end{kodekeluaran}

Penormalan telah Anda pelajari pada Bagian~5.2 untuk keperluan yang berbeda,
yaitu mencegah gerak diagonal yang terlalu cepat. Di sini penormalan dipakai
untuk keperluan yang lebih mendasar: memisahkan arah dari besar.

Tidak ada lagi nilai jarak pada perhitungan itu. Yang tersisa
hanyalah arah dan kecepatan yang Anda tetapkan. Inilah jawaban atas indikator
\ind{b} pada bab ini: besar kecepatan dibatasi oleh nilai yang Anda tentukan,
bukan oleh keadaan pemain.

\subsection{Perbandingan yang Dapat Diukur}

Perbedaan kedua cara itu dapat dibuktikan dengan mengukur kecepatan NPC pada
berbagai jarak. Program bab ini mencatat kecepatan NPC pertama setiap sepuluh
bingkai, beserta jaraknya saat itu. Tabel~\ref{tab:kecepatan-kejar} memuat
hasilnya pada satu komputer contoh.

\begin{table}[htbp]
\centering
\caption{Kecepatan terukur NPC pada berbagai jarak, dengan kecepatan ditetapkan 150 piksel per detik}
\label{tab:kecepatan-kejar}
{\setlength{\tabcolsep}{4pt}\footnotesize%
\begin{tabularx}{\textwidth}{@{}L{3.4cm} C{3.4cm} C{3.4cm} Y@{}}
\toprule
\headrow \textbf{Cara} & \textbf{Jarak saat diukur} & \textbf{Kecepatan terukur} & \textbf{Keterangan} \\
\midrule
Arah dinormalkan & 59--167 piksel & 150,0 piksel/detik & Tetap, tidak bergantung jarak. \\
\addlinespace
Vektor jarak dipakai langsung & 75--184 piksel & 75,2--184,0 piksel/detik &
Kecepatan sama persis dengan jaraknya. \\
\bottomrule
\end{tabularx}}
\tabsource{Hasil pengukuran penulis pada 400 bingkai dengan masukan otomatis.
Angka pada komputer Anda akan sedikit berbeda, tetapi polanya sama.}
\end{table}

Perhatikan baris kedua: kecepatannya sama persis dengan jaraknya. Itulah tanda
khas cara yang salah. Pada baris pertama, kecepatannya tetap walaupun jaraknya
berubah hampir tiga kali lipat.

\subsection{Mengapa Kecepatan yang Tetap Itu Penting}

Kecepatan yang bergantung pada jarak bukan sekadar kurang rapi; ia merusak
permainan dengan tiga cara.

\begin{enumerate}
  \item \textbf{Perilaku NPC tidak dapat diramalkan pemain.} Pemain tidak dapat
        belajar kapan ia aman dan kapan tidak, sebab kecepatan pengejar berubah
        menurut jaraknya. Pada permainan yang adil, pemain harus dapat
        memperkirakan akibat dari tindakannya.
  \item \textbf{Tingkat kesulitan berubah tanpa disengaja.} Ruang permainan
        yang luas membuat pemain lebih mudah lolos, sedangkan ruang yang sempit
        membuatnya lebih sulit. Tingkat kesulitan menjadi bergantung pada
        ukuran ruang, bukan pada rancangan Anda.
  \item \textbf{Penghitungan waktu menjadi keliru.} Apabila Anda memperkirakan
        NPC akan mencapai suatu tempat dalam dua detik, perkiraan itu hanya
        benar pada satu jarak tertentu.
\end{enumerate}

\begin{catatabox}
Sebagai pedoman yang berlaku umum: \textbf{arah dan besar kecepatan sebaiknya
ditentukan secara terpisah}. Arah berasal dari keadaan permainan, sedangkan
besarnya berasal dari rancangan Anda. Pedoman itu berlaku tidak hanya untuk
mengejar, melainkan juga untuk berpatroli, menghindar, dan seluruh perilaku
lain yang akan Anda susun pada bab-bab berikutnya.
\end{catatabox}
